% GENERATED FILE. Edit canonical lessons, then run npm run generate:books.
% canonical-source-hash: fnv1a64:2d856ade9892315d
% canonical-lessons: AR-C111-ishtaqa, AR-C111-tawaqqu, AR-C111-hamas, AR-C111-umniyya, AR-C111-tafdil

\chapter{To long for, Expectation, Enthusiasm, Wish, Preference}
\label{ch:ar-to-long-for-expectation-enthusiasm-wish-preference}
\hlchaptermodality{111}

\begin{chapteropening}
\textbf{By the end of this chapter:} \emph{I can say to long for, an expectation, enthusiasm, a wish and a preference.}

Complete the last lesson of chapter 111: I can say to long for, an expectation, enthusiasm, a wish and a preference.
\end{chapteropening}

\section[\texorpdfstring{\ar{اشتاق} (ishtāqa)}{اشتاق (ishtāqa)}]{\ar{اشتاق} (ishtāqa) — to long for}

\label{lesson:AR-C111-ishtaqa}



\begin{quote}
Before the new one: say the Arabic for to keep, then the Arabic for a talent.
\end{quote}

\subsection*{You'll want to know: \ar{اشتاق}}
\textbf{\ar{اشتاق}} — \emph{ishtāqa} — \textquotedblleft{}to long for\textquotedblright{}.

\begin{cousinweb}
\textbf{\ar{اشتاق}} (\emph{ishtāqa}) is built on the native Arabic root \emph{sh-w-q}, \textquotedblleft{}to long for\textquotedblright{} (Lane). \emph{Ishtāqa} is Form VIII, \textquotedblleft{}to feel longing\textquotedblright{}.
\end{cousinweb}

\begin{grammarlens}[title={What you've built}]
One more word for this chapter.
\end{grammarlens}

\subsection*{Your turn}
\begin{itemize}
\raggedright
  \item \emph{Say it:} \emph{ishtāqa}
  \item \emph{Say it:} \emph{ishtāqa}, once more
  \item \emph{Say it:} say \emph{ishtāqa}
  \item \emph{Recall:} say the Arabic for to keep, then the Arabic for a talent, then say \emph{ishtāqa} again
\end{itemize}

\begin{tcolorbox}[breakable,colback=teal!4,colframe=teal!35!black,title={Before you move on}]
What does \ar{اشتاق} mean? (\textquotedblleft{}To long for\textquotedblright{}.) Say it once more.
\end{tcolorbox}
\section[\texorpdfstring{\ar{توقع} (tawaqquʿ)}{توقع (tawaqquʿ)}]{\ar{توقع} (tawaqquʿ) — an expectation}

\label{lesson:AR-C111-tawaqqu}



\begin{quote}
Before the new one: say the Arabic for a talent, then the Arabic for to long for.

Say \textquotedblleft{}red\textquotedblright{} and \textquotedblleft{}blue,\textquotedblright{} masculine then feminine. (\emph{Aḥmar, ḥamrāʾ}; \emph{azraq, zarqāʾ}.)
\end{quote}

\subsection*{You'll want to know: \ar{توقع}}
\textbf{\ar{توقع}} — \emph{tawaqquʿ} — \textquotedblleft{}an expectation\textquotedblright{}.

\begin{cousinweb}
\textbf{\ar{توقع}} (\emph{tawaqquʿ}) is built on the native Arabic root \emph{w-q-ʿ}, \textquotedblleft{}to fall\textquotedblright{} (Lane). \emph{Tawaqquʿ} is the verbal noun of Form V, \textquotedblleft{}awaiting what will befall\textquotedblright{}.
\end{cousinweb}

\begin{grammarlens}[title={What you've built}]
One more word for this chapter.
\end{grammarlens}

\subsection*{Your turn}
\begin{itemize}
\raggedright
  \item \emph{Say it:} \emph{tawaqquʿ}
  \item \emph{Say it:} \emph{tawaqquʿ}, once more
  \item \emph{Say it:} say \emph{tawaqquʿ}
  \item \emph{Recall:} say the Arabic for a talent, then the Arabic for to long for, then say \emph{tawaqquʿ} again
\end{itemize}

\begin{tcolorbox}[breakable,colback=teal!4,colframe=teal!35!black,title={Before you move on}]
What does \ar{توقع} mean? (\textquotedblleft{}An expectation\textquotedblright{}.) Say it once more.
\end{tcolorbox}
\section[\texorpdfstring{\ar{حماس} (ḥamās)}{حماس (ḥamās)}]{\ar{حماس} (ḥamās) — enthusiasm}

\label{lesson:AR-C111-hamas}



\begin{quote}
Before the new one: say the Arabic for to long for, then the Arabic for an expectation.
\end{quote}

\subsection*{You'll want to know: \ar{حماس}}
\textbf{\ar{حماس}} — \emph{ḥamās} — \textquotedblleft{}enthusiasm\textquotedblright{}.

\begin{cousinweb}
\textbf{\ar{حماس}} (\emph{ḥamās}) is built on the native Arabic root \emph{ḥ-m-s}, \textquotedblleft{}to be hard, to be zealous\textquotedblright{} (Lane). Lane's \emph{ḥamāsa} is courage and zeal; \emph{ḥamās} is \textquotedblleft{}enthusiasm\textquotedblright{}.
\end{cousinweb}

\begin{grammarlens}[title={What you've built}]
One more word for this chapter.
\end{grammarlens}

\subsection*{Your turn}
\begin{itemize}
\raggedright
  \item \emph{Say it:} \emph{ḥamās}
  \item \emph{Say it:} \emph{ḥamās}, once more
  \item \emph{Say it:} say \emph{ḥamās}
  \item \emph{Recall:} say the Arabic for to long for, then the Arabic for an expectation, then say \emph{ḥamās} again
\end{itemize}

\begin{tcolorbox}[breakable,colback=teal!4,colframe=teal!35!black,title={Before you move on}]
What does \ar{حماس} mean? (\textquotedblleft{}Enthusiasm\textquotedblright{}.) Say it once more.
\end{tcolorbox}
\section[\texorpdfstring{\ar{أمنية} (umniyya)}{أمنية (umniyya)}]{\ar{أمنية} (umniyya) — a wish}

\label{lesson:AR-C111-umniyya}



\begin{quote}
Before the new one: say the Arabic for an expectation, then the Arabic for enthusiasm.
\end{quote}

\subsection*{You'll want to know: \ar{أمنية}}
\textbf{\ar{أمنية}} — \emph{umniyya} — \textquotedblleft{}a wish\textquotedblright{}.

\begin{cousinweb}
\textbf{\ar{أمنية}} (\emph{umniyya}) is built on the native Arabic root \emph{m-n-y}, \textquotedblleft{}to wish, to desire\textquotedblright{} (Lane). Lane's \emph{umniyya} is a thing wished for.
\end{cousinweb}

\begin{grammarlens}[title={What you've built}]
One more word for this chapter.
\end{grammarlens}

\subsection*{Your turn}
\begin{itemize}
\raggedright
  \item \emph{Say it:} \emph{umniyya}
  \item \emph{Say it:} \emph{umniyya}, once more
  \item \emph{Say it:} say \emph{umniyya}
  \item \emph{Recall:} say the Arabic for an expectation, then the Arabic for enthusiasm, then say \emph{umniyya} again
\end{itemize}

\begin{tcolorbox}[breakable,colback=teal!4,colframe=teal!35!black,title={Before you move on}]
What does \ar{أمنية} mean? (\textquotedblleft{}A wish\textquotedblright{}.) Say it once more.
\end{tcolorbox}
\section[\texorpdfstring{\ar{تفضيل} (tafḍīl)}{تفضيل (tafḍīl)}]{\ar{تفضيل} (tafḍīl) — a preference}

\label{lesson:AR-C111-tafdil}



\begin{quote}
Before the new one: say the Arabic for enthusiasm, then the Arabic for a wish.
\end{quote}

\subsection*{You'll want to know: \ar{تفضيل}}
\textbf{\ar{تفضيل}} — \emph{tafḍīl} — \textquotedblleft{}a preference\textquotedblright{}.

\begin{cousinweb}
\textbf{\ar{تفضيل}} (\emph{tafḍīl}) is built on the native Arabic root \emph{f-ḍ-l}, \textquotedblleft{}to exceed; excess, favour\textquotedblright{} (Lane). \emph{Tafḍīl} is the verbal noun of Form II, \textquotedblleft{}preferring\textquotedblright{}.
\end{cousinweb}

\begin{grammarlens}[title={What you've built}]
One more word for this chapter.
\end{grammarlens}

\subsection*{Your turn}
\begin{itemize}
\raggedright
  \item \emph{Say it:} \emph{tafḍīl}
  \item \emph{Say it:} \emph{tafḍīl}, once more
  \item \emph{Say it:} say \emph{tafḍīl}
  \item \emph{Recall:} say the Arabic for enthusiasm, then the Arabic for a wish, then say \emph{tafḍīl} again
\end{itemize}

\begin{tcolorbox}[breakable,colback=teal!4,colframe=teal!35!black,title={Before you move on}]
What does \ar{تفضيل} mean? (\textquotedblleft{}A preference\textquotedblright{}.) Say it once more.
\end{tcolorbox}
